% answers_short.tex —— 2011 年·浙江卷·理科 · 简短答案速查
% 由 tools/make_short_answers.py 自 answers.tex 的 \gksolhead/\gkanswer 提取，
% 并与 2011/2011_数学真题_浙江理_答案要点.md 逐题比对后生成，勿手改（改源 MD 后重跑）。

\begin{gknotice}
\textbf{使用说明}：本节只给最终答案便于快速核对，完整推导见解析版。答案为编写者独立推导，\textbf{非官方标准答案}。
\end{gknotice}

\gkpart{选择题}
\begin{center}\begin{tabular}{ccccccccccc}
\toprule
题号 & 1 & 2 & 3 & 4 & 5 & 6 & 7 & 8 & 9 & 10 \\
\midrule
答案 & B & B & C & D & B & C & A & C & B & D \\
\bottomrule
\end{tabular}\end{center}

\gkpart{填空题}
\begin{itemize}[leftmargin=2.6em]
  \item[\textbf{11.}] $0$。
  \item[\textbf{12.}] $5$。
  \item[\textbf{13.}] $2$。
  \item[\textbf{14.}] $\left[\dfrac\pi6,\dfrac{5\pi}6\right]$。
  \item[\textbf{15.}] $\dfrac53$。
  \item[\textbf{16.}] $\dfrac{2\sqrt{10}}5$。
  \item[\textbf{17.}] $(0,\pm1)$。
\end{itemize}

\gkpart{解答题}
\begin{itemize}[leftmargin=2.6em]
  \item[\textbf{18.}] (1) $a=1$，$c=\dfrac14$ 或 $a=\dfrac14$，$c=1$；(2) $\left(\dfrac{\sqrt6}2,\sqrt2\right)$
  \item[\textbf{19.}] (1) $a_n=na$，$S_n=\dfrac{a n(n+1)}2$；(2) $A_n<B_n$
  \item[\textbf{20.}] (1) 证明见解析；(2) 存在，$AM=3$
  \item[\textbf{21.}] (1) $\dfrac{17}4$；(2) $y=4\pm\dfrac{3\sqrt{115}}{115}x$
  \item[\textbf{22.}] (1) $a=\mathrm e$；(2) $\left[3\mathrm e-\dfrac{2\mathrm e}{\sqrt{1+\ln3}},\,3\mathrm e\right]$
\end{itemize}
